\documentclass[twocolumn]{article}

\usepackage[utf8]{inputenc}
\usepackage{listings}
\usepackage{cuted}

\usepackage[backend=biber, sorting=none]{biblatex}
\addbibresource{test.bib}
\usepackage{lipsum}
\usepackage{graphicx}
\usepackage{subcaption}

\usepackage[top=1in, bottom=1in, left=1.25in, right=1.25in]{geometry}

\usepackage[onehalfspacing]{setspace}
\usepackage[english]{babel}

\usepackage{amsmath, amssymb, amsthm}
\usepackage{enumerate, enumitem}
\usepackage{fancyhdr, graphicx, proof, comment, multicol}
\usepackage[none]{hyphenat}
\usepackage{dirtytalk}
\binoppenalty=\maxdimen
\relpenalty=\maxdimen
\usepackage{microtype}
%\usepackage{mathpazo}
\usepackage{bbm}
\usepackage{mdframed}
\usepackage{parskip}
\usepackage{pgfgantt}
\linespread{1.1}
\usepackage{graphicx}
\usepackage{subfig}
\usepackage{mathrsfs}
\usepackage{amsfonts}
\usepackage{amsmath}
\usepackage{minted}
\usepackage{amssymb}
\usepackage{algorithm}
\usepackage[noend]{algpseudocode}
\usepackage{mathtools}
\usepackage[colorlinks=false, linkcolor=black]{hyperref}
\usepackage[style=numeric,sorting=none,backend=biber]{biblatex}

\renewcommand*{\finalnamedelim}{\addcomma\space}
% \renewcommand*{\bibnumfmt}[1]{[#1]} % optional: adds square brackets

\title{Generating images from a subset of the ImageNet Dataset using Deep Convolutional Generative Adversarial
Networks (DCGANs)\\ \\Assigment 2}
\author{Aran Roig}
\date{\today}

%%%%%%%%%%%%%%%%%%%%%%%%%%%%%%%%%%%%%%%%%%%%%%%%%%%%%%%%%%%%%%%%%%%%%%%%%
% theorem environments
%%%%%%%%%%%%%%%%%%%%%%%%%%%%%%%%%%%%%%%%%%%%%%%%%%%%%%%%%%%%%%%%%%%%%%%%%

\newtheorem{defn0}{Definition}[section]
\newtheorem{prop0}[defn0]{Proposition}
\newtheorem{thm0}[defn0]{Theorem}
\newtheorem{lemma0}[defn0]{Lemma}
\newtheorem{corollary0}[defn0]{Corollary}
\newtheorem{example0}[defn0]{Example}
\newtheorem{remark0}[defn0]{Remark}
\newtheorem{conjecture0}[defn0]{Conjecture}

\newenvironment{definition}{ \begin{defn0}}{\end{defn0}}
\newenvironment{proposition}{\bigskip \begin{prop0}}{\end{prop0}}
\newenvironment{theorem}{\bigskip \begin{thm0}}{\end{thm0}}
\newenvironment{lemma}{\bigskip \begin{lemma0}}{\end{lemma0}}
\newenvironment{corollary}{\bigskip \begin{corollary0}}{\end{corollary0}}
\newenvironment{example}{ \begin{example0}\rm}{\end{example0}}
\newenvironment{remark}{ \begin{remark0}\rm}{\end{remark0}}
\newenvironment{conjecture}{\begin{conjecture0}}{\end{conjecture0}}

\newcommand{\defref}[1]{Definition~\ref{#1}}
\newcommand{\propref}[1]{Proposition~\ref{#1}}
\newcommand{\thmref}[1]{Theorem~\ref{#1}}
\newcommand{\lemref}[1]{Lemma~\ref{#1}}
\newcommand{\cororef}[1]{Corollary~\ref{#1}}
\newcommand{\exref}[1]{Example~\ref{#1}}
\newcommand{\secref}[1]{Section~\ref{#1}}
\newcommand{\remref}[1]{Remark~\ref{#1}}
\newcommand{\conjref}[1]{Conjecture~\ref{#1}}
\DeclarePairedDelimiterX{\norm}[1]{\lVert}{\rVert}{#1}

\begin{document}

\maketitle

\begin{strip}
\section*{Abstract}


Generative Adversarial Networks (GANs) have demonstrated remarkable capacity for
high-fidelity image synthesis, yet training stability and sample diversity remain
persistent challenges when scaling to large, semantically heterogeneous datasets. 

In this work, we investigate the use of conditional GANs for image generation on a
    curated subset of the ImageNet Large Scale Visual Recognition Challenge (ILSVRC) dataset \cite{russakovsky2014imagenet}, 
comprising 100 classes selected to span a range of visual complexity and intra-class variance.

    Our training pipeline incorporates 
    a two-timescale update rule (TTUR) to promote generator–discriminator equilibrium. 

    Qualitative analysis further confirms strong semantic coherence and inter-class disentanglement across 
    generated samples. These findings suggest that targeted subset training is a viable and efficient 
    strategy for domain-specific generative modeling, with implications for data augmentation, few-shot 
    learning, and controllable content generation.
\vspace{1cm}
\end{strip}
\section{Introduction}

Generative Adversarial Networks (GANs), introduced by Goodfellow et al. (2014) \cite{goodfellow2014generative} have become one of the most
widely studied frameworks for image generation. By putting a generator network against a discriminator
network in a minmax game, GANs are able to learn rich data representations and producing
visually convincing synthetic images.

Despite the proliferation of incresingly complex GAN architectures in recent literature, the core
framework remains a powerful and instructive baseline for image synthesis tasks.

In this work we apply a standard deep convolutional GAN (DCGAN) \cite{radford2015unsupervised} to a small subset of the
ImageNet dataset. A sample of this data can be seen on fig \ref{fig:sample} ImageNet is a large-scale visual database containing over fourteen million
labeled images spanning thousands of different categories, and the dataset has also served as a
foundational benchmark for image classification, generation and general computer vision research.


\begin{figure}[h]
    \centering
    \includegraphics[width=0.5\textwidth]{outputs/sample.png}
    \caption{Sample images of the used subsets of the ImageNet dataset}
    \label{fig:sample}
\end{figure}

Our model employs transposed convolutional layers in the generator to progressively upsample latent vectors
into full resolution image, while the discriminator uses strided convolutions to learn and
distinguish real samples from generated samples from the generator. Our goal is evaluating how well
conceptually a simple GAN architecture can capture the appearance of the structure of real-world categories.

\section{Literature Survey}

\subsection{Preceding Generative models}

GANs emerged in a moment where there were already pre-existing generative models, each with some significant limitations:

\textbf{Restricted Boltzmann Machines} (RBMs) and \textbf{Deep Boltzmann Machines} (DBMs) required
Markov chan Monte Carlo sampling, making the generation of synthetic data slow and appproximate.

\textbf{Variational Autoencoders (VAEs)} \cite{kingma2019introduction} offered an alternative deep generative framework with
principled latent variable inference but often produced blurry samples due to pixel-level reconstruction losses.

\textbf{Generative Stochastic Networks (GSNs)} \cite{alain2015gsns} extended denoising autoencoders but relied on complex approximate 
inference procedures

The GAN framework offered a compelling alternative by sidestepping explicit likelihood estimation, and
it proposed a trainable method that could produce sharper, more realistic samples.

\subsection{The original GAN model}

The concept of Generative Adversarial Networks was introduced by Ian J. Goodfellow, Jean Pouget-Abadie, Mehdi Mirza,
Bing Xu, David Warde-Farley, Sherjil Ozair, Aaron Courville and Yoshoua Beingo in 2014. The foundational paper titled
"Generative Adversarial Nets" \cite{goodfellow2014generative} proposed a framework for training generative models through an adversarial
process involving two models that competed between each other in a minimax game.

The two models of the GAN are trained simultaneously, each one having a different purpose:
\begin{itemize}
    \item A generative model $G$ captures the underlying data distribution of the training set. Its goal is to transform
        a noise vector $z$, called latent space, into synthetic samples.
    \item A discriminative model $D$ that estimates the probability that a given sample is coming from real data or if it is synthetic
        data generated from the generator $G$.
\end{itemize}

The adversarial objective can be formally expressed as:

$$\min_G \max_D V(D, G) = $$$$ = \mathbb{E}(\log{D(x)}) + \mathbb{E}(\log{(1 - D(G(z)))})$$

This adversarial min max game converges when the generator perfectly recovers the training data distribution
and the discriminator assigns a probability of $1/2$ to each given sample, so it is unable to discriminate between
real and synthetic data.

The original implementation of GAN networks used Multilayer Perceptrons (MLPs) as both the generator and the discriminator.
Experiments were conducted on three datasets:

\begin{itemize}
    \item The MNIST \cite{lecun1998gradient} dataset, which contains labeled images of handwritten digits
    \item The Toronto Face Dataset (TFD) \cite{susskind2010tfd}
    \item The CIFAR-10 natural image dataset \cite{krizhevsky2009learning}
\end{itemize}

The results of the original paper were promising but limited on image quality, particularly for complex
natural image domains. It demonstrated the potential of the framework while highlighting the need
for architectural improvements to scale to richer visual content.

\subsection{Early GAN variants and successors}

\textbf{Conditional GAN (cGAN)}: \cite{mirza2014conditional} Conditional GANs extend the original GAN framework by conditioning both the generator
and the discriminator on auxiliar information such as class labels. This allowed directed generation of specific
classes of data, converting the unsupervised formulation into one capable of target synthesis. cGANs demonstrated
that class dependant image generation could be achieved with a simple architectural extension.

\textbf{Laplacian Pyramid GANs (LAPGAN)}: \cite{denton2015deep} LAPGAN employed a cascade of convolutional networks with
a Laplacian pyramid framework, generating images in a coarse-to-fine fashion. While LAPGAN achieved visually
compelling results, the multi-network approach added more complexity. This framework demonstrated
that convolutional architectures were promising for adversarial image generation, directly motivating
the development of an unified deep convolutional architecture, that later would become DCGANs.

\subsection{Deep Convolutional Generative Adversarial Netwokrs (DCGAN)}

Despite the theoretical elegance of the original GAN, practical training proved to be difficult. The use of
fully connected architectures imposed strong independence assumptions, inappropiate for image data. The input data
of these models needed to have limited resolution because of its dense structure. The original GAN models also
produced noisy and incoherent outputs over complex training datasets.

The integration of CNNs, which had demonstrated unparalleled feature learning for images in supervised settings (eg. AlexNet,
VGG or GoogLeNet), into the GAN framework was a natural next step but required some careful architectural design.

DCGANs were first addressed in the paper "Unsupervised Representation Learning with Deep Convolutional Generative Adversarial
Networks". This paper has since accumulated over 14,000 citations, cementing its status as a landmark contribution in deep learning.

The central contribution of DCGAN is a set of empirically validated architectural guidelines that stabilize
GAN training when using deep convolutional networks:

\begin{itemize}
    \item Replace all pooling layers with strided convolutions in the discriminator and fractional-strided (transposed) convolutions
        in the generator model. This allows both networks to learn their own spatial downsampling and upsampling rather than
        relying on fixed pooling operations. Strided convolutions preserves more spatial information and allows learning spatial
        hirearchies.
    \item Use batch normalization in both the generator and the discriminator. Batch normalization stabilizes training by normalizing
        the input to each unit, mitigating issues with poor weight initialization and allowing gradients to flow more reliably through
        deeper networks.
    \item Remove fully connected hidden layers for deeper architectures. The original GAN models used dense layers to connect
        the latent vector to the convolutional layers. DCGANs instead connects directly from the latent vector to the first
        convolutional layer using a reshape operation. This enables a fully convolutional pipeline.
    \item Use ReLU activation in the generator for all layers except the output, which uses tanh. Use LeakyReLu (with a negative slope
        of $0.2$) in the discriminator for all layers, with a Sigmoid output for binary classification of real and fake data.
\end{itemize}

\subsection{Datasets in GAN and DCGAN Research}

The choice of picking a dataset is fundamental to GAN research, as it determines the complexity
of the distribution function to be learned, the resolution of the data and diversity of the input samples, and
the applicability of evaluation metrics. Different GAN architectures have been evaluated on progressively more
challenging datasets as the field has advanced.

\textbf{MNIST and CIFAR-10}. \cite{lecun1998gradient} \cite{krizhevsky2009learning} The MNIST dataset, comprising of $70 000$ grayscale images of handwritten digits
at $28 \times 28$ resolution was one of the first datasets used to evaluate the original GAN model in $2014$.
Its simplicity made it an ideal proof-of-concept benchmark. On the other hand, the CIFAR-10 dataset consisted
of $60 000$ color images across $10$ object classes at $32 \times 32$ resoluton. This dataset provided a more
challenging benchmark. Both datasets are continued to serve as standard evaluation platforms in deep learning research
due to their reference metrics and also wide adoption.

\textbf{LSUN - Large-Scale Scene Understanding} \cite{yu2015lsun} The LSUN dataset was one of the three primary datasets
used to evaluate DCGAN in the original paper. The bedroom susbet contains aproximately 3 million annotated
scene images, providing a large-scale training corpus. DCGAN trained on LSUN bedrooms was among the first GAN
models to generate photorealistic interior scenes at a $64 \times 64$ resolution, demonstrating the capability
of the convolutional adversarial approach.

\textbf{ImageNet-1k} \cite{russakovsky2014imagenet} ImageNet-1k contains approximately $1.2$ million images across $1000$ object categories. It was used
in the DCGAN paper to evaluate the quality of learned representations for supervised tasks. The discriminator's convolutional
features were used to train a regularized linear $L2-SVM$ classifier on the CIFAR-10 datset. This classifier achieved
an $82\%$ accuracy without any fine-tunning, demonstrating the transferability of representations learned using adversarial
training.

\textbf{Toronto Face Dataset (TFD)} \cite{susskind2010tfd} The Toronto Face Dataset was used in the original GAN paper as an early benchmark
for face generation. TFD contains approximately $100000$ grayscale facial images representing various expressions
and identitites. While it served as a foundational role in early GAN evaluations, it has been replaced by larger and
more diverse datasets such as CelebA and FFHQ for face generation research.

\section{Methodology}

We tested the DCGAN novel architecture against $100$ random classes of the ImageNet dataset. Our objective was to make
the generator learn the distribution of these $100$ classes, so we can generate synthetic data taken from these subset.
The implementation was made using the PyTorch python library. Additional libraries such as NumPy handled numerical operations,
and Matplotlib handled data visualization.

\subsection{Parameters}

Each layer of the generator has a convolutional transpose, followed by a batch normalization and a ReLU activation function.
It has $5$ of these layers:
\begin{itemize}
    \item $\text{ConvTranspose2d}(100 \rightarrow 512, k=4, s=1, p=0)$
    \item $\text{ConvTranspose2d}(512 \rightarrow 256, k=4, s=2, p=1)$
    \item $\text{ConvTranspose2d}(256 \rightarrow 128, k=4, s=2, p=1)$
    \item $\text{ConvTranspose2d}(128 \rightarrow 64, k=4, s=2, p=1)$
    \item $\text{ConvTranspose2d}(64 \rightarrow 3, k=4, s=2, p=1)$ (Tanh no batchnorm)
    \item \textbf{Output: }$3 \times 64 \times 64$
\end{itemize}
And the discriminator takes as an input the generated image, does some regular convolutions with batch normalization and a LeakyReLU with a negative slope of $0.2$, and finally
a sigmoid activation function:
\begin{itemize}
    \item $\text{Conv2d}(3 \rightarrow 64, k=4, s=2, p=1)$ (No batchnorm)
    \item $\text{Conv2d}(64 \rightarrow 128, k=4, s=2, p=1)$
    \item $\text{Conv2d}(128 \rightarrow 256, k=4, s=2, p=1)$
    \item $\text{Conv2d}(256 \rightarrow 512, k=4, s=2, p=1)$
    \item $\text{Conv2d}(512 \rightarrow 1, k=4, s=2, p=0)$
    \item \textbf{Output: }$1$
\end{itemize}

\section{Experimental results}

We measured the loss of the generator and the discriminator during each epoch of training. The results can be seen on fig. \ref{fig:gen_loss} and \ref{fig:disc_loss}.

An evolution of the model can be seen on the appendix on fig. \ref{fig:evolution}

\section{Conclusions}
As we can see, DCGANs provide a rich framework model for creating generative models of synthetic data that capture really well the underlying features of the training data,
and also can achieve good results on approximating the distribution function of the input dataset, as we have seen by training a test DCGAN model over a small slice of
the ImageNet dataset.

The implementation of the project (and the source of this document) can be found on \url{https://github.com/BinarySandia04/Assigment2ALLM}


\begin{figure}[h]
    \centering
    \includegraphics[width=0.5\textwidth]{outputs/gen_loss.png}
    \caption{The generator loss during training.}
    \label{fig:gen_loss}
\end{figure}

\begin{figure}[h]
    \centering
    \includegraphics[width=0.5\textwidth]{outputs/disc_loss.png}
    \caption{The discriminator loss during training.}
    \label{fig:disc_loss}
\end{figure}

\printbibliography

\appendix
\onecolumn
\section{Appendix}


\begin{figure*}[th]
    \centering
    
    % First row
    \begin{subfigure}{0.48\textwidth}
        \centering
        \includegraphics[width=\linewidth]{outputs/epoch_1.png}
    \end{subfigure}
    \hfill
    \begin{subfigure}{0.48\textwidth}
        \centering
        \includegraphics[width=\linewidth]{outputs/epoch_15.png}
    \end{subfigure}

    \vspace{0.5cm}

    % Second row
    \begin{subfigure}{0.5\textwidth}
        \centering
        \includegraphics[width=\linewidth]{outputs/epoch_30.png}
    \end{subfigure}
    \hfill
    \begin{subfigure}{0.48\textwidth}
        \centering
        \includegraphics[width=\linewidth]{outputs/epoch_50.png}
    \end{subfigure}

    \caption{Evolution of the output of the generator over the same noise sample on different epochs.
    One can see that on later epochs the quality of the synthetic images improves.}
    \label{fig:evolution}
\end{figure*}


\end{document}
